\section{Autenticazione passwordless con FIDO2}\label{sec:fido2}

Pazienti e medici accedono senza password, con una \emph{passkey} FIDO2. Questa sezione descrive registrazione, login, sessioni e revoca, e le scelte che li rendono resistenti agli attacchi del modello di minaccia.

\subsection{WebAuthn in breve}
FIDO2 è formato da WebAuthn\cite{webauthn}, l'interfaccia tra browser e server, e da CTAP2, il protocollo tra browser e autenticatore. I componenti sono tre:
\begin{itemize}
  \item l'\textbf{autenticatore}, \emph{platform} (Windows Hello, il sensore del telefono) o \emph{roaming} (una chiavetta USB), che genera una coppia di chiavi per ogni sito e non esporta mai la chiave privata;
  \item il \textbf{client}, cioè il browser, che inoltra le richieste e aggiunge l'\emph{origin} della pagina;
  \item la \textbf{Relying Party}, il nostro backend, che conserva solo la chiave pubblica e verifica le firme.
\end{itemize}
Registrazione e login sono schemi \emph{challenge-response} a chiave pubblica\cite{kaufman}: il server invia un valore casuale, l'autenticatore lo firma dopo aver verificato l'utente con impronta, volto o PIN, e il server controlla la firma. Sul server non c'è alcun segreto da rubare, e la firma è legata all'origin: una pagina clone, anche identica nell'aspetto, ha un origin diverso e la sua risposta viene rifiutata. È questo che rende WebAuthn resistente al phishing.

Il backend usa Node.js~24, TypeScript, Express~5 e la libreria \texttt{@simplewebauthn/server}~14 per generare le opzioni e verificare le risposte; credenziali e sessioni sono conservate in MongoDB.

\subsection{Registrazione}
Alla richiesta delle opzioni il server risponde con una challenge casuale di 32 byte, valida 5 minuti, e con questi parametri:
\begin{itemize}
  \item \texttt{userVerification: required}: l'autenticatore deve verificare l'utente, quindi la passkey vale come due fattori (possesso del dispositivo e biometria o PIN);
  \item \texttt{residentKey: required}: la passkey è \emph{individuabile}, cioè l'autenticatore conserva anche l'identificativo dell'utente, così al login è il dispositivo a proporla;
  \item un \texttt{userID} casuale, senza dati personali;
  \item \texttt{attestation: none}: il server non chiede il certificato del produttore dell'autenticatore (vedi i rischi residui).
\end{itemize}
Alla verifica la challenge viene consumata prima di ogni controllo, così una risposta non può essere ritentata né riproposta; la libreria controlla tipo, challenge, origin, \texttt{rpID}, flag di presenza e verifica dell'utente e algoritmo. Un indice univoco nel database impedisce di associare la stessa credenziale a due utenti.

\paragraph{Chi può registrare una passkey.} La verifica crittografica dice che la passkey è autentica, non che chi la registra abbia il diritto di farlo per quell'utente: senza un controllo, chiunque potrebbe aggiungere la propria passkey all'account di un altro e accedere al suo posto (\emph{account takeover}). La prima passkey si registra quindi solo con un \textbf{codice d'invito}, che la struttura consegna dopo aver identificato la persona (\emph{identity proofing}): 128 bit casuali, legato a uno username, monouso, valido 24 ore, conservato sul server solo come hash. Le passkey successive le aggiunge solo l'utente già autenticato (listato~\ref{lst:autorizza}). L'invito viene consumato con un'operazione atomica e solo a registrazione riuscita: due registrazioni contemporanee con lo stesso codice non possono riuscire entrambe.

\begin{lstlisting}[language=TypeScript,caption={Autorizzazione alla registrazione (\texttt{controllers/registration.ts}, estratto).},label={lst:autorizza}]
// Può registrare una passkey per "username":
// - l'utente stesso, già autenticato (aggiunta di un dispositivo);
// - chi presenta un invito valido per quello username.
async function autorizza(store: Store, req: Request, username: string) {
  if ((await sessionUser(store, req)) === username) return { tipo: "sessione" };
  const invitation = req.body?.invitation;
  if (typeof invitation === "string" && invitation.length > 0) {
    const codeHash = hashToken(invitation);
    if (await store.hasInvitation(codeHash, username)) return { tipo: "invito", codeHash };
  }
  return null;
}
\end{lstlisting}

\subsection{Login}
Il login segue lo stesso schema; il server verifica la firma con la chiave pubblica registrata, oltre a challenge, origin, \texttt{rpID} e verifica dell'utente. Quattro scelte lo proteggono da attacchi meno ovvi.
\begin{itemize}
  \item \textbf{Opzioni uguali per tutti.} Le opzioni di login non elencano le credenziali dell'utente (\texttt{allowCredentials} vuoto): sono identiche per qualunque username, esistente o no, e non permettono di scoprire chi è registrato (\emph{user enumeration}). Funziona perché le passkey sono individuabili.
  \item \textbf{Challenge indicizzate per valore.} La challenge non è legata allo username: alla verifica il server la legge dai dati firmati dal browser e la consuma con un'operazione atomica (listato~\ref{lst:login}). Se fosse legata allo username, chiunque potrebbe chiedere opzioni di login a nome di un medico e sostituirgli la challenge, impedendogli di entrare.
  \item \textbf{Contatore anti-riuso.} Ogni autenticatore incrementa un contatore (\emph{signCount}) a ogni firma, e un valore che non cresce rivela una chiave clonata. Il contatore viene aggiornato con un \emph{compare-and-set}: la scrittura riesce solo se nel database c'è ancora il valore letto prima della verifica, così se un clone e l'originale accedono nello stesso istante ne passa uno solo.
  \item \textbf{Errori uniformi.} Utente inesistente, credenziale sconosciuta, firma non valida e contatore errato producono la stessa risposta; il motivo è scritto solo nel log del server.
\end{itemize}

\begin{lstlisting}[language=TypeScript,caption={Challenge e contatore al login (\texttt{controllers/authentication.ts}, estratto).},label={lst:login}]
// La challenge si legge dai dati firmati dal browser e si consuma in modo atomico.
const challenge = challengeOf(response);
if (!challenge || !(await store.takeLoginChallenge(challenge))) {
  res.status(400).json({ error: "challenge assente o scaduta" });
  return;
}
// ... verifica della firma ...
// Il nuovo signCount si salva solo se nessun altro login lo ha cambiato nel frattempo.
const updated = await store.updateCounter(username, credential.id,
  expectedCounter, verification.authenticationInfo.newCounter);
if (!updated) {
  res.status(401).json({ error: "autenticazione non valida" });
  return;
}
\end{lstlisting}

\subsection{Algoritmi e crittografia post-quantistica}\label{sec:fido2-pq}
Il server offre, in ordine di preferenza, gli algoritmi ML-DSA-44, EdDSA, ES256 e RS256, e l'autenticatore sceglie il primo che supporta. ML-DSA è lo schema di firma post-quantistico basato su reticoli standardizzato dal NIST in FIPS~204\cite{fips204}, resistente all'algoritmo di Shor che romperebbe ECDSA, EdDSA e RSA. La codifica delle chiavi ML-DSA per WebAuthn è definita in RFC~9964\cite{rfc9964}, e i primi gestori di passkey generano già credenziali ML-DSA\cite{1password-pq}; la maggior parte dei dispositivi, Windows Hello compreso, usa ancora ES256. Il sistema accetta entrambe senza modifiche al codice (\emph{crypto-agility}), in linea con il piano di transizione del NIST, che depreca le firme classiche dopo il 2030\cite{nistir8547}. Per l'autenticazione il rischio quantistico riguarda il futuro: una firma intercettata oggi non si può riusare, perché legata a una challenge monouso. Il rischio \emph{harvest now, decrypt later} riguarda invece la riservatezza dei referti, protetti con ML-KEM (sezione~\ref{sec:dati}).

\subsection{Sessioni e revoca}\label{sec:fido2-revoca}
Dopo il login il server rilascia un token di sessione in un cookie, con le scelte della tabella~\ref{tab:sessioni}, che seguono le raccomandazioni OWASP\cite{owasp-session}.

\begin{table}[H]
\centering\footnotesize
\renewcommand{\arraystretch}{1.15}
\begin{tabular}{@{}p{.3\textwidth}p{.62\textwidth}@{}}
\toprule
Scelta & Motivazione\\
\midrule
Token opaco di 256 bit casuali & Impossibile da indovinare; a differenza di un JWT si revoca subito.\\
Sul server solo l'hash SHA-256 & Chi legge l'archivio delle sessioni non trova token utilizzabili.\\
Cookie \texttt{HttpOnly}, \texttt{SameSite=Strict}, \texttt{Secure} in produzione & Non leggibile da script iniettati (XSS), non inviato da altri siti (CSRF), solo su HTTPS.\\
Token nuovo a ogni login & Protezione dalla \emph{session fixation}.\\
Scadenza dopo 15 minuti di inattività e comunque dopo 8 ore & Una postazione lasciata aperta non resta autenticata; un token rubato non vale per sempre.\\
Middleware unico \texttt{requireSession} & Un solo punto di controllo per tutte le rotte protette (\emph{complete mediation}).\\
\bottomrule
\end{tabular}
\caption{Scelte di sicurezza per le sessioni.}\label{tab:sessioni}
\end{table}

Per un dispositivo perso, l'utente può revocare una propria passkey: il server chiude tutte le sue sessioni, compresa quella eventualmente aperta dal dispositivo perso. L'ultima passkey non si può revocare, e la rimozione è atomica, così due revoche contemporanee non lasciano l'utente senza passkey. L'utente può anche uscire da tutti i dispositivi. Chi ha perso tutti i dispositivi si rivolge alla struttura, che dopo averlo identificato revoca d'ufficio tutte le sue passkey (\texttt{npm run revoca}) e gli consegna un nuovo invito.

\subsection{Persistenza su MongoDB}\label{sec:fido2-mongo}
MongoDB gira in un contenitore Docker raggiungibile solo dal computer stesso (\texttt{127.0.0.1}), con autenticazione obbligatoria: il backend usa un utente con il solo ruolo \texttt{readWrite} sul database dell'applicazione, e se venisse compromesso non potrebbe leggere altri database né creare utenti. Alcune proprietà di sicurezza sono garantite dal database stesso: un indice univoco impedisce che la stessa credenziale appartenga a due utenti, indici TTL cancellano da soli challenge, sessioni e inviti scaduti (minimizzazione dei dati), e ogni operazione che deve avvenire una volta sola, come consumare una challenge o un invito, è eseguita in un unico passaggio atomico (\texttt{findOneAndDelete}).

\subsection{Verifica}\label{sec:fido2-test}
I test usano un \textbf{autenticatore software} che genera chiavi ES256 o ML-DSA-44 e costruisce risposte WebAuthn come un dispositivo vero. A differenza di un dispositivo vero, può essere manomesso di proposito (origin falso, contatore riportato indietro), e permette così di riprodurre gli attacchi in modo ripetibile. La tabella~\ref{tab:test-fido2} riporta i principali, con il riferimento alle minacce della sezione~\ref{sec:minacce}.

\begin{table}[H]
\centering\footnotesize
\renewcommand{\arraystretch}{1.15}
\begin{tabular}{@{}p{.36\textwidth}p{.07\textwidth}p{.47\textwidth}@{}}
\toprule
Attacco simulato & Rif. & Esito\\
\midrule
Replay di una risposta di login & M2 & Rifiutato: challenge già consumata\\
Phishing in registrazione e al login & M1 & Rifiutato: origin non atteso\\
Autenticatore clonato, anche usato in parallelo all'originale & M3 & Rifiutato il contatore non crescente; dei due login contemporanei ne passa uno\\
Registrazione senza verifica dell'utente & M1 & Rifiutata\\
Passkey dell'attaccante su un account altrui; inviti assenti, usati, scaduti o di altri & M12 & Rifiutata (403)\\
Richieste di login a nome di un altro utente & M10 & Il login della vittima riesce comunque\\
Confronto delle opzioni di login & M13 & Risposte indistinguibili\\
Token inventato, cookie dopo il logout, sessione scaduta o oltre 8 ore & M6, M17 & Rifiutati (401)\\
Sessione di un dispositivo poi revocato & M17 & Chiusa; la passkey revocata non entra più\\
Firma ML-DSA-44 prodotta da un'altra chiave & M1 & Rifiutata\\
\bottomrule
\end{tabular}
\caption{Principali test di sicurezza dell'autenticazione.}\label{tab:test-fido2}
\end{table}

I test di integrazione su MongoDB verificano le stesse proprietà quando dipendono dal database: il contatore salvato blocca un clone anche dopo un riavvio, nel database c'è solo l'hash dei token, e challenge, inviti e contatori non possono essere usati due volte da richieste parallele. Il flusso completo è stato provato anche in un browser reale con una passkey di Windows Hello sbloccata con l'impronta: nel database risultano solo identificativo, chiave pubblica, contatore e trasporto \texttt{internal}, che indica una chiave conservata nel chip di sicurezza del computer.

\subsection{Rischi residui}
\begin{itemize}
  \item \textbf{Passkey sincronizzate.} Quelle salvate nel cloud spesso inviano sempre contatore zero, e il controllo anti-clonazione non si applica: per i professionisti sanitari sono preferibili passkey legate al dispositivo.
  \item \textbf{Attestazione.} Senza attestazione il server non sa quale autenticatore abbia creato la passkey; per i medici si potrebbero ammettere solo autenticatori hardware certificati, a costo di più complessità e di qualche informazione in più sul dispositivo.
  \item \textbf{Operazioni sensibili senza nuova verifica.} Revocare una passkey richiede solo una sessione valida: una nuova verifica con la passkey (\emph{step-up}) renderebbe innocua una postazione lasciata aperta.
  \item \textbf{Inviti e recupero.} Un codice d'invito, finché non è usato, equivale a una password temporanea e va consegnato su un canale sicuro; la sicurezza del recupero coincide con quella dell'identificazione presso la struttura.
  \item \textbf{HTTPS.} In produzione il servizio va esposto solo su HTTPS, con \texttt{rpID} uguale al dominio reale.
\end{itemize}
